
\begin{exercise}[subtitle={Uncorrelated vs independent \Coffeecup\Coffeecup}]
   Let \(X\) and \(Y\) be two random variables in \(\real\). Prove that
   the following statements are equivalent:
   \begin{enumerate}[label=(\roman*),noitemsep]
        \item\label{it: independent} \(X\) and \(Y\) are independent, that is
        \[
            \prob{X \in A, Y \in B} = \prob{X \in A}\prob{Y \in B} \qquad \forall A, B
        \]
        \item\label{it: measurable functions} For all measurable functions \(f\) and \(g\) such that \(f(X)\)
        and \(g(X)\) are integrable we have
        \[
            \E{f(X)g(Y)} = \E{f(X)}\E{g(Y)}.
        \]
        \item\label{it: bounded measurable functions} For all \underline{bounded}, measurable functions \(f\) and \(g\)
        \[
            \E{f(X)g(Y)} = \E{f(X)}\E{g(Y)}.
        \]
   \end{enumerate}
   \textbf{Hint: } Use Theorem 102 from Ivan's lecture notes.
    
   \noindent Further prove the following statements
    \begin{enumerate}[label=(\alph*),noitemsep]
        \item\label{it: idep => uncor}
        If \(X\) and \(Y\) are independent, they are uncorrelated, that is
        \[
            \cov{X}{Y} = 0.
        \]

        \item\label{it: uncor not indep}
        Let \(\prob{X=-1} = \prob{X=0} = \prob{X=1} = \frac13\) and \(Y =
        X^2\). Prove that \(X\) and \(Y\) are uncorrelated but not independent.

        \item\label{it: indep stable under transformation}
        If \(X\) and \(Y\) are independent, then \(h_1(X)\) and \(h_2(Y)\) are independent
        for all measurable functions \(h_1\colon \real \to \real\) and \(h_2\colon \real \to \real\).

        \item\label{it: all transformations uncorrelated implies independence}
        If \(f(X)\) and \(g(Y)\) are uncorrelated for all bounded
        measurable functions \(f\) and \(g\), then \(X\) and \(Y\) are
        independent.
    \end{enumerate}

    \begin{remark}[Interpretation]
        What this essentially means is the following:
        \begin{enumerate}
            \item Uncorrelated is weaker than independence.
            \item \(f(X)\) and \(f(Y)\) are independent if \(X\) and \(Y\)
            are independent, that is independence is \emph{stable under transformation}.

            \item In contrast, \(X\) and \(Y\) being uncorrelated is \emph{not} stable under transformation.
            In fact, if it were true that \(f(X)\) and \(g(Y)\) is uncorrelated
            for all bounded measurable functions \(f\) and \(g\), then \(X\) and
            \(Y\) would be independent by \ref{it: all transformations
            uncorrelated implies independence}. Independence therefore gains the
            interpretation of ``all transformations are uncorrelated''.
        \end{enumerate}
    \end{remark}
\end{exercise}

\begin{solution}
\begin{enumerate}[wide=0pt]
    \item[\ref{it: independent} \(\Leftrightarrow\) \ref{it: measurable functions} \(\Leftrightarrow\) \ref{it: bounded measurable functions}:]

    We have \ref{it: independent} \(\Rightarrow\) \ref{it: measurable functions} by Theorem 102 from Ivan's lecture.
    The implication \ref{it: measurable functions} \(\Rightarrow\) \ref{it: bounded measurable functions} is trivial.
    For the implication \ref{it: bounded measurable functions} \(\Rightarrow\) \ref{it: independent} observe that for any \(A, B\) we have
    that \(\ind{A}\) and \(\ind{B}\) are bounded measurable functions.

    \item[\ref{it: idep => uncor}]
    We have that
    \begin{align*}
        \cov{X}{Y}
        &= \E{(X-\E{X})(Y-\E{Y})}
        \\
        &= \E{XY} - \E{X}\E{Y}
        \overset{\text{indep.}}=
        \E{X}\E{Y} - \E{X}\E{Y} = 0.
    \end{align*}

    \item[\ref{it: uncor not indep}]
    Since \(\E{X} = 0\) we have that
    \begin{align*}
        \cov{X}{Y}
        &= \E{XY} - \E{X}\E{Y}
        \\
        &= \E{X^3} - \E{X}\E{X^2}
        = \E{X} - \E{X}\E{X^2}
        = 0.
    \end{align*}
    However, 
    \[
        \prob{X=0, Y=1} =0 \neq \tfrac13 \tfrac23 = \prob{X=0}\prob{Y=1}.
    \]

    \item[\ref{it: indep stable under transformation}]
    Since \(f\circ h_1\) and \(g \circ h_2\) are measurable for all
    measurable \(f\colon \real \to \real\) and \(g\colon \real \to \real\), we have by \ref{it: measurable
    functions} that
    \begin{align*}
        \E{f(h_1(X))g(h_2(Y))}
        &= \E{(f\circ h_1)(X)(g\circ h_2)(Y)}
        \\
        &= \E{f(h_1(X))}\E{g(h_2(Y))}.
    \end{align*}
    This proves independence by \ref{it: measurable functions}. 

    \item[\ref{it: all transformations uncorrelated implies independence}]
    We have
    \[
        0 = \cov{f(X)}{g(Y)} = \E{f(X),g(Y)} - \E{f(X)}\E{g(Y)}
    \]
    and consequently \(\E{f(X)g(Y)} = \E{f(X)}\E{g(Y)}\) for all
    bounded measurable functions \(f\) and \(g\).  This implies
    independence by \ref{it: bounded measurable functions}.
\end{enumerate}
\end{solution}
